% =============================================================================
% DRAFT -- not part of the report. Nothing under report/ inputs this file.
%
% Written outside report/ deliberately. The report describes agentic_full at
% exact match 0.432 and 0.224, and that remains true of the system it
% documents: its eighteen evaluated runs are untouched and its tables still
% regenerate byte-identically from main. Whether any of this chapter enters
% the report is the owner's decision, not the executor's.
%
% To include it, move it to report/chapters/ch6_improvement.tex, add
% \input{chapters/ch6_improvement} to report/main.tex after ch4, and
% regenerate the three tables it inputs:
%
%   python -c "from agentic_ir.eval.tables import write_v2_tables as w; print(*w(), sep=chr(10))"
%
% Every number below cites the run that produced it. Run identifiers refer to
% directories under results/runs_v2 (the loop's) or results/runs (the report's).
% =============================================================================

\chapter{Improving the System, and What It Cost}
\label{ch:improvement}

\section{What Was Attempted}
\label{sec:improvement-attempted}

Chapter~\ref{ch:implementation} reports a system that loses to its own
baselines on answer quality. That is an honest result and it is not a
satisfying one, so a bounded improvement loop was run against it: five
hypotheses, fixed in advance, each measured on the frozen 250-question
evaluation slice, each kept or dropped on a paired comparison rather than on a
point estimate. The loop was constrained to stop on stated conditions rather
than when a target was reached, so that a null could end it as legitimately as
a success.

The constraints were those of the project throughout. One consumer GPU, one
local model, no paid API. The nine evaluated configurations of
Chapter~\ref{ch:implementation} were not to change: every modification ships
behind a configuration flag defaulting to the evaluated behaviour, so that
\texttt{agentic\_full} remains the system the report describes and its tables
continue to regenerate byte-identically from the runs on disk. That property
is checked rather than asserted --- the four evaluated agentic configurations
resolve to identical configuration trees before and after the loop's changes.

One hypothesis survived.

\section{The Kept Change: Reasoning on the Synthesis Call}
\label{sec:improvement-kept}

\texttt{qwen3:8b} is a hybrid reasoning model. Throughout
Chapter~\ref{ch:implementation} its reasoning mode is suppressed --- by the
\texttt{think} flag, by a \texttt{/no\_think} suffix in the system prompt, and
by stripping the block from responses --- because at roughly 30--45 tokens per
second an unbudgeted reasoning block costs ten to twenty seconds per call and
the pipeline makes several calls per question.

The kept configuration, \texttt{agentic\_v2}, re-enables it for exactly one
agent. The synthesiser reasons; the Planner, Retrieval Agent, Knowledge Graph
Navigator and Verifier do not. This is not a global switch: \texttt{llm.think}
is one flag passed down verbatim by every caller, which cannot express
``reasoning on the synthesiser only'', so a per-agent routing list was added
beside it. The configuration also raises the completion budget from 1{,}024 to
5{,}120 tokens, and that is part of the change rather than incidental to it: at
1{,}024 the reasoning block consumes the entire budget and the synthesiser
returns no answer at all.

% GENERATED by src/agentic_ir/eval/tables.py::write_v2_tables -- do not edit by hand.
% Source runs: results/runs_v2 (the improvement loop). `cli tables` neither
% writes nor reads these -- discover_runs cannot see that directory, which is
% what keeps the report's nine-system grid exactly nine systems.
% Requires \usepackage{booktabs} and \usepackage{graphicx}.

\begin{table}[htbp]
  \centering
  \small
  \setlength{\tabcolsep}{4pt}
  \caption{The improvement loop's kept configuration against the reported system, its own base, and the strongest answer baseline. Paired on question id over the same frozen 250-question slice. Five of the six intervals exclude zero; the exception is \texttt{self\_ask} on HotpotQA, where the kept configuration ties a single-model baseline that has no agents in it.}
  \label{tab:v2-results}
  \resizebox{\ifdim\width>\textwidth\textwidth\else\width\fi}{!}{%
  \begin{tabular}{lrrrlrr}
    \toprule
    System & $n$ & EM & F1 & $\Delta$EM vs. kept [95\% CI] & $p$ & gained\,/\,lost \\
    \midrule
    \multicolumn{7}{@{}l}{\textit{HotpotQA}} \\
    \textbf{agentic\_v2} (kept) & 250 & \textbf{0.524} & \textbf{0.618} & -- & -- & -- \\
    \texttt{agentic\_full} & 250 & 0.432 & 0.510 & +0.092 [+0.036, +0.144] & $<$0.001 & 36\,/\,13 \\
    \texttt{agentic\_no\_planner} & 250 & 0.456 & 0.543 & +0.068 [+0.024, +0.112] & 0.002 & 25\,/\,8 \\
    \texttt{self\_ask} & 250 & 0.504 & 0.603 & +0.020 [-0.032, +0.072] & 0.499 & 25\,/\,20 \\
    \addlinespace
    \multicolumn{7}{@{}l}{\textit{2WikiMultihopQA}} \\
    \textbf{agentic\_v2} (kept) & 250 & \textbf{0.456} & \textbf{0.511} & -- & -- & -- \\
    \texttt{agentic\_full} & 250 & 0.224 & 0.267 & +0.232 [+0.168, +0.296] & $<$0.001 & 67\,/\,9 \\
    \texttt{agentic\_no\_planner} & 250 & 0.324 & 0.388 & +0.132 [+0.080, +0.188] & $<$0.001 & 42\,/\,9 \\
    \texttt{self\_ask} & 250 & 0.264 & 0.310 & +0.192 [+0.128, +0.256] & $<$0.001 & 64\,/\,16 \\
    \bottomrule
  \end{tabular}}
\end{table}


Table~\ref{tab:v2-results} gives the outcome, paired on question identifier
over the same 250 questions, with 2{,}000 bootstrap resamples. Against the
system this report describes, exact match rises from 0.432 to 0.524 on HotpotQA
and from 0.224 to 0.456 on 2WikiMultihopQA. Both intervals exclude zero. The
runs are \texttt{L1c\_think\_hotpotqa} and \texttt{L1c\_think\_twowiki}.

The comparison that matters most for the architectural claim is the one that
does not clear zero. On HotpotQA the kept configuration gains $+0.020$ over
\texttt{self\_ask}, with an interval of $[-0.032, +0.072]$ and $p = 0.499$: it
ties a single-model baseline that contains no agents at all. On
2WikiMultihopQA the same comparison gives $+0.192$. Stated plainly, the
architecture earns its place on the dataset built to require multiple hops and
does not earn it on the dataset that does not. This is the same conclusion
Chapter~\ref{ch:implementation} reaches about \texttt{agentic\_full}, reached
again at a higher absolute score, and the improvement does not overturn it.

% GENERATED by src/agentic_ir/eval/tables.py::write_v2_tables -- do not edit by hand.
% Source runs: results/runs_v2 (the improvement loop). `cli tables` neither
% writes nor reads these -- discover_runs cannot see that directory, which is
% what keeps the report's nine-system grid exactly nine systems.
% Requires \usepackage{booktabs} and \usepackage{graphicx}.

\begin{table}[htbp]
  \centering
  \small
  \setlength{\tabcolsep}{4pt}
  \caption{Cost per question. The kept configuration answers with roughly a third of the model calls and a faster median, and pays for it in the tail: a synthesis call that reasons may run to the full 5{,}120-token completion budget.}
  \label{tab:v2-cost}
  \resizebox{\ifdim\width>\textwidth\textwidth\else\width\fi}{!}{%
  \begin{tabular}{llrrr}
    \toprule
    Dataset & System & Median (s) & p90 (s) & LLM calls \\
    \midrule
    HotpotQA & \textbf{agentic\_v2} (kept) & \textbf{18.7} & 115.6 & \textbf{1.37} \\
     & \texttt{agentic\_full} & 28.3 & 55.4 & 4.26 \\
    2WikiMultihopQA & \textbf{agentic\_v2} (kept) & \textbf{32.2} & 144.7 & \textbf{1.21} \\
     & \texttt{agentic\_full} & 35.4 & 69.2 & 3.80 \\
    \bottomrule
  \end{tabular}}
\end{table}


Table~\ref{tab:v2-cost} gives the cost. The kept configuration answers with
roughly a third of the model calls of \texttt{agentic\_full} and a faster
median, because the one-node plan it builds on issues a single sub-query rather
than a decomposition. It pays in the tail: a synthesis call that reasons may run
to the full 5{,}120-token budget, and the p90 roughly doubles. For a system
whose per-question latency was already the dominant practical constraint, that
trade is worth stating in both directions rather than reporting the median
alone.

\section{What the Change Costs That the Scores Do Not Show}
\label{sec:improvement-caveat}

Enabling reasoning introduces a failure mode that suppressing it does not have.

% GENERATED by src/agentic_ir/eval/tables.py::write_v2_tables -- do not edit by hand.
% Source runs: results/runs_v2 (the improvement loop). `cli tables` neither
% writes nor reads these -- discover_runs cannot see that directory, which is
% what keeps the report's nine-system grid exactly nine systems.
% Requires \usepackage{booktabs} and \usepackage{graphicx}.

\begin{table}[htbp]
  \centering
  \small
  \setlength{\tabcolsep}{4pt}
  \caption{What reasoning on the synthesis call costs. The model spends its entire generation budget inside the thinking channel and never opens the content channel (\texttt{done\_reason} \texttt{length} at the \texttt{num\_predict} cap), and the extractive fallback answers instead. Every one of these losses is already inside the exact-match figures of Table~\ref{tab:v2-results}.}
  \label{tab:v2-failures}
  \resizebox{\ifdim\width>\textwidth\textwidth\else\width\fi}{!}{%
  \begin{tabular}{lrrrr}
    \toprule
    Condition & $n$ & Synthesis calls & Questions losing one & Questions retrying \\
    \midrule
    HotpotQA, reasoning off & 250 & 250 & 0 (0.0\%) & 0 (0.0\%) \\
    2WikiMultihopQA, reasoning off & 250 & 250 & 0 (0.0\%) & 0 (0.0\%) \\
    HotpotQA, reasoning on & 250 & 251 & 2 (0.8\%) & 26 (10.4\%) \\
    2WikiMultihopQA, reasoning on & 250 & 251 & 8 (3.2\%) & 49 (19.6\%) \\
    \bottomrule
  \end{tabular}}
\end{table}


Table~\ref{tab:v2-failures} counts it. Across 500 questions with reasoning
suppressed, no synthesis call failed. Across the same 500 with reasoning
enabled, ten failed outright and a further seventy-five required a retry that
the suppressed condition never needed. The mechanism is exact and was confirmed
by instrumenting the client: the call ends with \texttt{done\_reason} equal to
\texttt{length} at an \texttt{eval\_count} of exactly 5{,}120, the configured
cap. The model spends its entire generation budget inside the thinking channel
and never opens the content channel. Ollama returns reasoning out of band, so
the visible reply is genuinely empty rather than truncated mid-sentence, and the
orchestrator's extractive fallback rule then supplies a confident wrong answer
in its place.

Two things follow, and both matter for how Table~\ref{tab:v2-results} should be
read. First, those ten losses are already inside the reported figures: each
scores zero exact match in the runs that produced 0.524 and 0.456, so the gains
are net of them rather than reduced by them. Second, the obvious repair does not
work. Re-issuing each of the ten failed calls with reasoning disabled recovers
the content channel in ten cases out of ten, on the first attempt, and produces
zero correct answers. These are questions the system was going to lose in any
case; the extractive fallback that currently answers them is no worse than the
model would have been. The seventy-five retries are therefore a latency cost and
not an accuracy one.

This failure was invisible for two rounds of the loop. The trace schema declared
a \texttt{truncated} field and no code ever assigned it, so it read false on
every call ever recorded, including these ten. The loop concluded twice, in
writing, that the empty completions were \emph{not} the completion cap. Both
conclusions were wrong, and the field that would have settled the question was
sitting in the record. It is now populated from \texttt{done\_reason}.

\section{A Methodological Finding: Seed Pinning Is Not Reproducibility}
\label{sec:improvement-reproducibility}

Every measurement in this report is made at \texttt{temperature} zero with
\texttt{seed} 42 and \texttt{PYTHONHASHSEED} pinned. That is not sufficient for
determinism on this stack, and the loop established why.

Issuing an identical prompt repeatedly, in one process, at temperature zero,
with the same seed and the same model, produces different completions. The
prompt token count reported for the same first-attempt prompt varied across
repetitions --- 1{,}218, 2{,}977 and 3{,}671 --- because Ollama's key-value
cache reuses the shared prefix of a preceding request. A different prefix reuse
gives a different batch shape, a different batch shape gives different
floating-point arithmetic, and the sampling that follows diverges even when it
is nominally greedy.

This supplies a mechanism for an observation made earlier in the project and
left unexplained: re-running byte-identical code over the same slice moves 6--14
percent of answers. It is not a defect in the harness and no amount of seed
discipline removes it. Its consequence is quantitative and applies to every
comparison in this report. At $n = 250$, an effect smaller than roughly three
exact-match points cannot be distinguished from the harness itself. Differences
above that threshold --- the $+0.092$ and $+0.232$ of
Table~\ref{tab:v2-results}, and the ablation gaps of
Chapter~\ref{ch:implementation} --- survive it. Differences below it, including
several reported there as unresolved, are correctly described as unresolved and
would remain so under any number of repetitions of the same design.

\section{What Was Dropped, and Why That Is Worth Reporting}
\label{sec:improvement-dropped}

Four of the five hypotheses were dropped. Two were closed without spending GPU
time at all, which is the loop's most useful property rather than an
embarrassment.

\paragraph{A larger local model.} The clearest hypothesis was that every result
in this report is conditional on a 4-bit 8B model and that a larger one would
lift it. It does not fit. \texttt{qwen3:14b} runs at 11.1 tokens per second on
this hardware against the 24.5 the evaluation budget requires, with 58 percent
of the model resident on the GPU and the remainder on the CPU. A sub-hypothesis
--- that shrinking the context window would free memory for more GPU layers ---
was falsified precisely: freeing 0.62\,GiB of key-value cache bought
approximately zero additional layers, because the runtime pins the resident
fraction regardless. \texttt{gemma3:12b}, the fallback, reports no reasoning
capability at all and so cannot reproduce the one change that worked.

\paragraph{A wider evidence budget.} The plan held that the synthesiser's
20-item evidence cap was the binding constraint on answer quality, citing a
gain from 152 to 155 complete gold evidence pools when raised to 30. An offline
replay before any GPU time showed the premise was false for the system being
improved. The one-node plan issues a single sub-query returning three documents
and a median of ten passage sentences; only three questions of 250 exceed twenty
items. Raising the cap adds fourteen passage sentences across the entire slice
and moves complete pools from 142 to 143. The 152-to-155 figure described the
five-sub-query pool of \texttt{agentic\_full} --- a number carried across from a
different system, and the offline gate is what caught it.

\paragraph{Anchoring with a graph quota.} The same replay closed the next
hypothesis without a separate experiment. Knowledge-graph rows score between
0.2255 and 1.0000 in the evidence ranking while passage sentences are
normalised into 0.0333 to 0.1000, so every graph triple outranks every passage.
Widening the budget fills the new slots with triples placed above the sentences
that carry the gold answer. The hypothesis had proposed a reserved graph quota
as the fix; the measurement says the ranking is the defect, which a quota does
not address.

\paragraph{Self-consistency voting.} Sampling the synthesiser three times and
taking the majority answer was costed at 6.5 GPU hours for HotpotQA, within
budget, and the failure mode it targets was counted before spending them. Of
119 wrong answers, nineteen have the shape voting suppresses: the gold string
present in the model's own answer sentence but a different span in the answer
field. Fifteen of those nineteen, however, are surface-form artifacts ---
\texttt{ten} against \texttt{10}, \texttt{650} against \texttt{650 locations},
\texttt{Los Alamos} against \texttt{Los Alamos Laboratory} --- which are
deterministic formatting choices with no variance for voting to suppress. Three
draws would agree with one another and still disagree with the reference. Four
answers are genuinely different spans, giving a ceiling of $+0.016$ exact match
if voting corrected every one of them. Against a harness noise floor of roughly
three points, that effect is not resolvable at $n = 250$ even if it is entirely
real, so the run's interval would necessarily include zero and the hypothesis
was dropped on the count rather than on the run.

\section{Conclusion}
\label{sec:improvement-conclusion}

The loop ended on its stated stop condition --- two consecutive drops after the
performance floor had been reached --- with the backlog fully decided. It moved
exact match from 0.432 to 0.524 on HotpotQA and from 0.224 to 0.456 on
2WikiMultihopQA through a single change: allowing one agent of five to reason.
It did not close the gap to \texttt{self\_ask} on HotpotQA, and the honest
reading of Table~\ref{tab:v2-results} is that the agentic architecture is
vindicated on multi-hop questions and not on this benchmark as a whole.

Two of the four drops were decided offline, at no computational cost, because
their premises were checked against the system actually being improved rather
than assumed from the plan that proposed them. In one case the premise had been
imported from a different configuration and was simply wrong for this one. That
is the argument for measuring a hypothesis' precondition before measuring the
hypothesis, and it saved more time here than the one accepted change did.

The finding most likely to outlive this project is
Section~\ref{sec:improvement-reproducibility}: on a local inference server with
prefix caching, temperature zero and a fixed seed do not make a system
reproducible, and any evaluation of this size carries an irreducible noise floor
that no seed discipline removes. Reported effects should be checked against that
floor before they are believed --- in this report as much as in any other.
